\subsection{Orthogonality. Direct sums of bilinear or sesquilinear maps.}

In this section, A and B denote rings, E a left A-module, F a right (resp. left) B-module, G an (A, B)-bimodule, and Φ a bilinear (resp. sesquilinear for a given antiautomorphism J of B) map from E × F to G.

DEFINITION 6. --- Two elements \(x \in E\) and \(y \in F\) are said to be orthogonal with respect to \(\Phi\) if \(\Phi(x, y) = 0\) . Two subsets \(E' \subset E\) and \(F' \subset F\) are said to be orthogonal if, for any \(x \in E'\) and \(y \in F'\) , \(x\) and \(y\) are orthogonal. The set of elements of E (resp. F) orthogonal to a given submodule N of F (resp. M of E) is a submodule of E (resp. F), called the totally orthogonal (or simply orthogonal) submodule to N (resp. M), and denoted \(N^0\) (resp. \(M^0\) ).

Let H and H' be two submodules of E or of F. One has \(\mathrm{H}\subset(\mathrm{H}^{0})^{0}\) (denoted \(H^{00}\) ); if \(H\subset H'\) , one has \(H^{\prime0}\supset H^{0}\) . It follows that one has \(H^{0}\supset(H^{00})^{0}\) and \(H^{0}\subset(H^{0})^{00}\) ; by setting

\[
\mathrm{H} ^ {0 0 0} = (\mathrm{H} ^ {0 0}) ^ {0} = (\mathrm{H} ^ {0}) ^ {0 0} = ((\mathrm{H} ^ {0}) ^ {0}) ^ {0},
\]

we therefore have \(\mathbf{H}^0 = \mathbf{H}^{000}\) .

For the map \(\Phi\) to be degenerate (no. 1, def. 3) it is necessary and sufficient that at least one of the two submodules E \(^{0}\) , F \(^{0}\) be \(\neq \{0\}\) . It is clear that \(\Phi(x, y)\) does not change when one adds to \(x\) (resp. \(y\) )

an element of F \(^{0}\) (resp. E \(^{0}\) ), and Φ therefore defines, by passing to the quotient, a bilinear (or sesquilinear) map on (E/F \(^{0}\) ) × (F/E \(^{0}\) ); this is visibly nondegenerate; it is called the nondegenerate bilinear (or sesquilinear) map associated with Φ.

Let \((\mathrm{E}_{i})_{i\in\mathrm{I}}\) be a family of left A-modules, \((\mathrm{F}_{i})_{i\in\mathrm{I}}\) a family of right B-modules (resp. left B-modules), \(\Phi_{i}\) a bilinear map (resp. sesquilinear on the right for J) from \(E_{i}\times F_{i}\) into G. Let E (resp. F) denote the direct sum module of the \(E_{i}\) (resp. \(F_{i}\) ). We see immediately that the map \(\Phi\) of \(E\times F\) into G defined by

\[
\Phi ((x _ {i}), (y _ {i})) = \sum_ {i} \Phi_ {i} (x _ {i}, y _ {i}) \quad (x _ {i} \in \mathrm{E} _ {i}, y _ {i} \in \mathrm{F} _ {i})\tag{10}
\]

(a sum that makes sense since its terms are zero except for finitely many of them) is bilinear (resp. sesquilinear on the right for J). It is called the direct sum of the maps \(\Phi_{i}\) . It is clear that \(\mathbf{E}_{i}\) is orthogonal to \(\mathbf{F}_{j}\) with respect to \(\Phi\) for \(i \neq j\) . Conversely, let \(\Phi\) be a bilinear or sesquilinear map from \(\mathbf{E} \times \mathbf{F}\) into G, and suppose that E is the direct sum of submodules \((\mathbf{E}_{i})_{i \in I}\) and F the direct sum of submodules \((\mathbf{F}_{i})_{i \in I}\) such that \(\mathbf{E}_{i}\) is orthogonal to \(\mathbf{F}_{j}\) for \(i \neq j\) ; then \(\Phi\) is the direct sum of its restrictions to the products \(\mathbf{E}_{i} \times \mathbf{F}_{i} (i \in I)\) .

For \(\Phi\) to be nondegenerate, it is necessary and sufficient that each of the \(\Phi_{i}\) be so; under these conditions, the submodule orthogonal to \(\mathbf{E}_{i}\) is \(\sum_{j\neq i}\mathbf{F}_{j}\) .

